\section{局部系统的叠}\srcpage{23}

在 Langlands 对偶一侧，主角是局部系统的叠。本节讨论它在不同情形下的定义。

此处最熟悉的情形是 Betti 情形。于是 $X$ 是 $k=\C$ 上的光滑射影曲线，等价地说，是紧 Riemann 面。为简单起见，我们排除 $\mathbb{P}^1$ 的情形，并假设 $X$ 的亏格 $g\ge 1$。那么作为拓扑空间，$X$ 是其基本群 $\pi_1=\pi_1(X)$ 的 $K(\pi_1,1)$，而该群是有限表现群 —— 事实上，它有著名的表现
\[
\pi_1=\langle a_1,b_1,a_2,b_2,\dots,a_g,b_g\mid [a_1,b_1][a_2,b_2]\cdots[a_g,b_g]=1\rangle,
\]
有 $2g$ 个生成元与一个关系。

我们可以用任意特征 $0$ 的代数闭域 $e$ 作系数；读者可假设 $e=\C$，但 $e=\Ql$ 也将是相关的（不过无论如何，它也同构于 $\C$）。在此情形，$X$ 上的 $\Gd$-局部系统无非由同态
\[
\pi_1\to\Gd(e)
\]
给出，至多差 $\Gd(e)$-共轭。更确切地说，有仿射概形 $\Hom(\pi_1,\Gd)$ 参数化同态 $\pi_1\to\Gd$，而局部系统的叠是所得的商叠
\[
\LocSys^{\Betti}_{\Gd}=\Hom(\pi_1,\Gd)/\Gd.
\]

\srcpage{24}具体地，利用 $\pi_1$ 的表现，有
\[
\Hom(\pi_1,\Gd)=\Gd^{2g}\times_{\Gd}*,
\]
它是映射
\[
\Gd^{2g}\to\Gd\colon (a_1,b_1,\dots,a_g,b_g)\mapsto [a_1,b_1]\cdots[a_g,b_g]
\]
在 $1\in\Gd$ 处的纤维。

我们指出，一般情形下 —— 例如当 $\Gd$ 是环面时 —— 这个映射不是平坦的，而纤维积
\[
\Hom(\pi_1,\Gd)=\Gd^{2g}\times_{\Gd}*
\]
必须在导出意义下取，作为导出仿射概形，其函数环是导出张量积
\[
\cO(\Gd^{2g})\otimes^L_{\cO(\Gd)}e.
\]

若 $\Gd$ 半单且 $g\ge 2$，这个问题不会出现，例如参见 Aizenbud 与 Avni（2016）。一般地，$\Hom(\pi_1,\Gd)$ 只有温和的奇点：正如该表现所示，它是完全交。

因此，在 Betti 情形，局部系统的叠是极为基本且被充分研究的对象：至少当 $\Gd=\mathrm{GL}_n$ 时，它不过是视为叠的特征簇。

若 $e=\C$，则 Riemann--Hilbert 对应给出局部系统与带平坦联络的向量丛之间的翻译。这就把我们带到局部系统叠的 de Rham 体现。它有相同的 $\C$-点集（更确切地说是群胚），但代数结构不同。更一般地，若 $k$ 是任一特征 $0$ 的（代数闭）域，可取 $e=k$，并定义代数叠
\[
\LocSys^{\dR}_{\Gd},
\]
其 $A$-值点（对 $k$-代数 $A$ 而言）是 $X\times_{\Spec(k)}\Spec(A)$ 上带 $X$-方向平坦联络的 $\Gd$-丛构成的群胚。等价地，它们是 $X_{\dR}\times_{\Spec(k)}\Spec(A)$ 上的 $\Gd$-丛，其中 $X_{\dR}$ 是 $X$ 的 Simpson de Rham 叠。回忆：这是 $X$ 按对角线的形式完备化所作的商 $X_{\dR}=X/(X\subset X\times X)^{\wedge}$，而 $X_{\dR}$ 上的拟凝聚层恰是 $X$ 上的 D-模。

遗忘联络给出代数叠的映射
\[
\LocSys^{\dR}_{\Gd}\to\Bun_{\Gd},
\]
它在 $\Gd$-丛 $E$ 处的纤维分类 $E$ 上的联络。$E$ 上联络的集合构成以仿射空间 $\Gamma(X,E\times^{\Gd}\check{\gc}\otimes\OmX)$ 为模型的拟挠子。同样，一切都必须在导出代数几何的世界中解释，故 $A$ 是 $D_{\ge 0}(k)$ 中的代数，等等。在这一解释下，$\LocSys^{\dR}_{\Gd}$ 仍是局部完全交 —— $\Bun_{\Gd}$ 光滑，而到它的投影是局部完全交。关键在于复形 $\Gamma(X,E\times^{\Gd}\check{\gc}\otimes\OmX)$ 集中在上同调次数 $0$ 与 $1$。

最后考虑 étale 情形；此时 $e=\Ql$，且可允许特征不同于 $\ell$ 的任意代数闭域 $k$。在此情形，有投射有限群 $\pi_1^{\et}=\pi_1^{\et}(X)$，当 $k=\C$ 时它是\srcpage{25}上文所述离散群 $\pi_1$ 的投射有限完备化。若 $k$ 有正特征 $p$，则通过选取整环 $\cO_C$ 上的光滑射影曲线 $X_{\cO_C}$（$\cO_C$ 为某剩余域为 $k$ 的完备代数闭 $p$ 进域 $C$ 的整数环），并使 $X_{\cO_C}$ 形变 $X=X_k$，得到特殊化映射
\[
\pi_1^{\et}(X_C)\to\pi_1^{\et}(X),
\]
它是满射，把 $\pi_1^{\et}(X)$ 写成离散群 $\pi_1$ 的投射有限完备化的商。

理想情形下，人们会考虑 $\pi_1^{\et}$ 的连续表示的某个模空间；然而这样的模空间充其量只能生活在刚性解析几何的领域。就我们目前的目的而言，我们希望得到通常（导出）代数几何中的模叠。

粗略地说，Arinkin、Gaitsgory、Kazhdan、Raskin、Rozenblyum 与 Varshavsky（2020b）采取的办法，是把 $\pi_1^{\et}$ 换成它在 $\Ql$ 上的 pro-代数完备化，并考虑其表示。在 Tannakian 视角下，这意味着考虑全子范畴
\[
\Lisse(X)\subset D^b_c(X,\Ql),
\]
即 lisse $\Ql$-层的有界复形，然后与其 Ind-范畴 $\Ind\Lisse(X)$ 打交道。于是对任一导出 $\Ql$-代数 $A$，可以考虑左 $t$-正合的 $\otimes$-函子
\[
\Rep(\Gd)\to\Ind\Lisse(X)\otimes_{\Ql}A;
\]
这定义了叠 $\LocSys^{\et}_X$（在文献中记作 $\LocSys^{\restr}_X$）。\footnote{对 $X=\mathbb{P}^1$，需要作轻微的修改。}

为理解这个模叠的性质，我们指出，在 Betti 情形也可以这样做。此时令
\[
\LocSys^{\coarse}_{\Gd}=\Hom(\pi_1,\Gd)/\!/\Gd=\Spec(\cO(\LocSys^{\Betti}_{\Gd}))
\]
为粗略模空间。那么 Betti 叠的限制版本是
\[
\LocSys^{\Betti}_{\Gd}\times_{\LocSys^{\coarse}_{\Gd}}\bigsqcup_{x\in\LocSys^{\coarse}_{\Gd}(\C)}(\LocSys^{\coarse}_{\Gd})^{\wedge}_{x},
\]
即在粗略模空间层面只允许无穷小变化。模叠 $\LocSys^{\et}_{\Gd}$ 则是连通分支的并，对应于那些使相应的 $\Gd$-表示 $\pi_1\to\Gd(e)=\Gd(\Ql)$ 连续延拓到 $\pi_1^{\et}$ 的 $x$。
